\documentclass[9pt,twocolumn]{extarticle}
\usepackage[margin=0.67in]{geometry}
\usepackage{graphicx,booktabs,amsmath,hyperref,microtype,xcolor}
\usepackage[font=small,skip=3pt]{caption}
\setlength{\columnsep}{0.24in}
\setlength{\parindent}{0pt}\setlength{\parskip}{3pt}
\hypersetup{colorlinks=true,urlcolor=blue,citecolor=blue}
\begin{document}
\twocolumn[{\centering\Large\bfseries Detecting Political Stance in Social Media Posts Using Machine Learning\par\vspace{5pt}\normalsize Team: \textbf{Arek Bedros Kalayjian, Rami Hassan}\\ Emails: \textbf{abk25@mail.aub.edu, rnh33@mail.aub.edu}\\\vspace{7pt}}]
\section*{1. Problem and motivation}

Can a simple text classifier identify whether a tweet is \textsc{favor}, \textsc{against}, or \textsc{none} toward a specified topic, and does doubling the amount of training data meaningfully improve performance? Our specific hypothesis is that a classifier trained on 100\% rather than 50\% of available training examples will improve $F_{\mathrm{avg}}=(F_1^{\rm favor}+F_1^{\rm against})/2$ by \textbf{at least 0.02} on the same held-out evaluation set, averaged across five randomly selected half-size training samples.

This matters to social scientists tracking public debate and analysts studying how opinions vary across topics. Predicting stance, unlike ordinary sentiment, requires interpreting a statement relative to a target even when the tweet never names that target \cite{mohammad2016}.
\section*{2. Impact}
A successful result would give researchers evidence about whether gathering more labeled tweets is worthwhile for a lightweight stance detector. If the hypothesis fails, that is still informative: data quantity alone may not overcome sparse vocabulary, topic shifts, irony, and ambiguous wording. Automated predictions should not be treated as statements about individual political beliefs. Topic imbalance, non-representative Twitter sampling, noisy labels and asymmetric errors (especially falsely labeling \textsc{favor} as \textsc{against}) can distort conclusions. Results will be reported by topic and class, not only as an aggregate.
\section*{3. Related work}
\textbf{Mohammad et al.} introduced SemEval-2016 Task 6 and its target-aware stance labels; the official evaluation is the mean of FAVOR and AGAINST F1, excluding NONE. Their majority-class baseline scored \textbf{0.6522}; the strongest submitted Task A system reached \textbf{0.6782} \cite{mohammad2016}.

\textbf{Zarrella and Marsh} used transfer learning and achieved the top submitted Task A score, demonstrating that learned representations can capture patterns beyond straightforward word matching \cite{zarrella2016}. \textbf{Tutek et al.} combined diverse classifiers using a genetic-algorithm ensemble and placed near the top of the shared task \cite{tutek2016}. \textbf{Patra et al.} investigated support vector machines for tweets, motivating a comparison with other linear text classifiers \cite{patra2016}. \textbf{Vijayaraghavan et al.} explored word- and character-level convolutional networks with data augmentation, reporting $F_{\mathrm{avg}}\approx0.635$ \cite{vijayaraghavan2016}. \textbf{Augenstein et al.} investigated representations for stance generalization across targets, illustrating the importance of domain and topic differences \cite{augenstein2016}.

Unlike these systems, this project prioritizes an interpretable \textbf{training-set-size experiment}, with multinomial logistic regression implemented from scratch, not an attempt to beat the shared-task leaderboard.
\section*{4. Data}
We will use the publicly accessible \textbf{SemEval-2016 Task 6, Task A} corpus, curated by Mohammad et al. \cite{mohammad2016}, with a downloadable copy at \url{https://github.com/emsrc/SemEval2016_T6_Stance_Detection} \cite{repository}. It consists of English tweets collected for five targets: Atheism, Climate Change, Feminist Movement, Hillary Clinton, and Legalization of Abortion. Data were annotated for stance toward the target. Reuse will be limited to course research with source attribution; the mirror does not clearly state a separate dataset license, so redistribution permissions must be confirmed against the original task terms.

Each row includes an identifier (nonpredictive), target (categorical text), tweet (free text), and stance (three-class categorical response). The mirror's training file contains \textbf{2,814} labeled tweets (the shared-task paper reports 2,914, which likely also counts the trial tweets released for training) and the official test set has \textbf{1,249} examples. There are no physical units; bag-of-words will yield a sparse numeric vocabulary learned only from training data. Class balance differs significantly between targets, creating a risk of topic-specific majority-label shortcuts.
\begin{table}[h]
\centering\small
\caption{Corpus summary; train count measured locally by \texttt{baseline.py}, test count from \cite{mohammad2016}.}
\begin{tabular}{lr}\toprule Property & Value\\\midrule Task & A (five known targets)\\ Train examples & 2,814\\ Official test examples & 1,249\\ Input fields & ID, target, text\\ Classes & FAVOR, AGAINST, NONE\\ Evaluation metric & Mean F1 of FAVOR/AGAINST\\\bottomrule\end{tabular}
\end{table}
\begin{figure}[h]
\centering\includegraphics[width=\columnwidth]{class_distribution.png}
\caption{Training-set class counts by target, generated by \texttt{baseline.py}.}
\end{figure}
Data checks will count duplicates, empty tweets, missing target/stance fields and per-topic class frequencies. Preprocessing will preserve stance-bearing negations and hashtags, normalize case/URLs, and vectorize tweet text and its target. Some sarcasm, implicit stance, abbreviations, and sampling biases will remain.
\section*{5. Proposed models and methods}
\textbf{M0: Majority-class dummy.} Predict the most frequent class in the training subset. It tests how far imbalance alone takes us.

\textbf{M1: From-scratch softmax regression.} Encode lowercased word unigrams and the target with sparse bag-of-words counts. Train a three-output linear softmax model with cross-entropy, gradient descent and optional $L_2$ regularization. This suits small, high-dimensional text and lets us inspect feature weights.

\textbf{M2: Bigram softmax regression.} Add word bigrams to capture short phrases such as ``not support.'' Hold all other settings fixed where possible to isolate the feature contribution. A scikit-learn implementation will serve only as a verification reference \cite{pedregosa2011}; the final core learner will be our own implementation.

\textbf{Preliminary benchmark (published, not locally run):} On the official Task A test data, majority-class $F_{\mathrm{avg}}=0.6522$, SVM with unigrams $=0.6331$, and SVM with n-grams $=0.6898$ \cite{mohammad2016}. \textbf{Own training/validation results (80/20 split of the training file, scikit-learn reference models):} majority class val $F_{\mathrm{avg}}=0.3233$ (acc 0.4778); unigram logistic regression val $F_{\mathrm{avg}}=0.6085$ (train 0.9792, acc 0.6075); unigram+bigram val $F_{\mathrm{avg}}=0.5993$ (train 0.9987, acc 0.6021). The large train/validation gap shows overfitting that $L_2$ tuning should address. Our majority baseline predicts one global class (AGAINST), so it scores lower than the published per-target majority baseline.
\section*{6. Evaluation}
The primary metric is the official $F_{\mathrm{avg}}$, to avoid allowing a strong NONE score to obscure errors in the stance classes. We will also report accuracy, per-class precision/recall/F1, confusion matrices, and per-target F1. First, stratify the official training data into 80\% development training and 20\% validation by target and class. Keep the official test set untouched until the final assessment. Use validation data to select preprocessing and regularization settings.

For the main experiment, hold the architecture and hyperparameters fixed and sample five 50\% subsets of the available development-training pool without replacement (five fixed random seeds). Compare their performance on the official test set to the corresponding 100\% model; report the mean F1 difference, range and standard deviation. To avoid choosing settings using test performance, finish model selection using validation first. The hypothesis is supported only if the full-model improvement over the \emph{mean} of the five half-data models is at least 0.02. Additional comparisons: M0 vs M1 vs M2, with/without bigrams, and F1 across topics. A small qualitative error review will examine sarcasm/implicit stance.
\section*{7. Resources and code}
The dataset contains only a few thousand short documents. NumPy, pandas, Matplotlib and a sparse feature representation will suffice on a regular laptop; neither GPU nor AUB HPC is needed. A reproducible download/EDA/baseline script is included with this proposal. Public repository: \url{https://github.com/RamiHassan33/Stance-Classifier}. The repository will contain \texttt{baseline.py}, \texttt{README.md}, dependency instructions, exploratory figures, model code, and run logs; each team member will commit throughout the project.
\section*{8. Timeline}
\begin{table}[h]\centering\small\caption{Project milestones.}\begin{tabular}{p{0.22\columnwidth}p{0.7\columnwidth}}\toprule When & Deliverable\\\midrule Oct 20--25 & Data checks, repository, baseline and plots\\ Oct 26--Nov 1 & From-scratch softmax, gradient checks\\ Nov 2--8 & Bigram model, validation tuning\\ Nov 9--16 & Five half-data trials, class/topic analyses\\ \textbf{Nov 17} & \textbf{Progress report due}\\ Nov 18--29 & Error analysis, ablations, reproducibility\\ Dec 4 & Poster session\\ Dec 10 & Final report\\\bottomrule\end{tabular}\end{table}
\clearpage
\begin{thebibliography}{9}\small
\bibitem{mohammad2016} Mohammad et al. (2016). SemEval-2016 Task 6: Detecting Stance in Tweets. \emph{SemEval}, 31--41. \url{https://aclanthology.org/S16-1003/}.
\bibitem{zarrella2016} Zarrella and Marsh (2016). MITRE at SemEval-2016 Task 6: Transfer Learning for Stance Detection. \url{https://aclanthology.org/S16-1074/}.
\bibitem{tutek2016} Tutek et al. (2016). TakeLab at SemEval-2016 Task 6: Stance Classification in Tweets Using a Genetic Algorithm Based Ensemble. \url{https://aclanthology.org/S16-1075/}.
\bibitem{patra2016} Patra et al. (2016). JU\_NLP at SemEval-2016 Task 6: Detecting Stance in Tweets Using Support Vector Machines. \url{https://aclanthology.org/S16-1071/}.
\bibitem{vijayaraghavan2016} Vijayaraghavan et al. (2016). DeepStance at SemEval-2016 Task 6: Detecting Stance in Tweets Using Character and Word-Level CNNs. \url{https://aclanthology.org/S16-1067/}.
\bibitem{augenstein2016} Augenstein et al. (2016). USFD at SemEval-2016 Task 6: Any-Target Stance Detection on Twitter with Autoencoders. \url{https://aclanthology.org/S16-1063/}.
\bibitem{repository} emsrc (2016). SemEval2016 Task 6 Stance Detection (data/code mirror). \url{https://github.com/emsrc/SemEval2016_T6_Stance_Detection}.
\bibitem{pedregosa2011} Pedregosa et al. (2011). Scikit-learn: Machine Learning in Python. \emph{JMLR} 12, 2825--2830.
\end{thebibliography}
\end{document}
